\section*{Capítulo \ref{ch:g7:geological-time} --- As idades da Terra}

\begin{solution}{exo:g7:geological-time:1}
(a) \omterm{def:g7:geological-time:unconformity}{discordância}; (b) inclusão; (c) superposição; (d) sucessão dos
\omterm{def:g7:sediments-and-fossils:fossil}{fósseis}; (e) relações de corte.
\end{solution}

\begin{solution}{exo:g7:geological-time:2}
(a) O arenito vermelho, embaixo, é o mais velho; o arenito amarelo, em
cima, é o mais novo. (b) O \omterm{def:g7:geological-time:superposition}{princípio da superposição}. (c) As camadas não
foram viradas desde que se formaram.
\end{solution}

\begin{solution}{exo:g7:geological-time:3}
(a) Paleozoico, Mesozoico, Cenozoico. (b) No Cenozoico. (c) No
Mesozoico. (d) O Jurássico vem logo depois do Triássico; o Cretáceo vem
logo antes do Paleógeno.
\end{solution}

\begin{solution}{exo:g7:geological-time:4}
(a) Falso: a \omterm{def:g7:geological-time:relative-dating}{datação relativa} só põe as \omterm{def:g3:stones-pebbles-sand:rock}{rochas} em ordem; uma idade em
anos é uma \omterm{def:g7:geological-time:relative-dating}{idade absoluta}. (b) Verdadeiro. (c) Falso: os \omterm{def:g3:stones-pebbles-sand:pebble}{seixos} são mais
velhos que o conglomerado que os contém (\omterm{def:g7:geological-time:inclusion}{princípio da inclusão}).
(d) Verdadeiro (\omterm{def:g7:geological-time:fossil-succession}{princípio da sucessão dos fósseis}).
\end{solution}

\begin{solution}{exo:g7:geological-time:5}
(a) Mais novo: ele corta a camada 3, e uma \omterm{def:g3:stones-pebbles-sand:rock}{rocha} que corta outra é mais
nova (relações de corte). (b) Mais velho: a camada 4 fica por cima dele
sem ser cortada. (c) A camada 1, a camada 2, a camada 3, o granito, a
camada 4. (d) As \omterm{def:g3:stones-pebbles-sand:rock}{rochas} acima dele foram desgastadas pela \omterm{def:g5:water-shapes-the-land:erosion}{erosão} até o
granito chegar à superfície; a camada 4 foi então depositada sobre essa
superfície de \omterm{def:g5:water-shapes-the-land:erosion}{erosão}, uma \omterm{def:g7:geological-time:unconformity}{discordância}.
\end{solution}

\begin{solution}{exo:g7:geological-time:6}
(a) $538.8 - 251.9 = 286.9$ milhões de anos. (b) $251.9 - 66 = 185.9$
milhões de anos. (c) 66 milhões de anos. (d) O Paleozoico é a mais longa,
o Cenozoico a mais curta (até agora).
\end{solution}

\begin{solution}{exo:g7:geological-time:7}
(a) $4540 - 538.8 = 4001.2$ milhões de anos. (b) $4001.2 \div 4540
\approx 0.881$: cerca de \textbf{88~\%}. (c) Antes do Paleozoico, os
animais com concha ainda não eram comuns, e os seres vivos moles quase
nunca viram \omterm{def:g7:sediments-and-fossils:fossil}{fósseis}; as conchas e os ossos, que se fossilizam bem, só são
comuns a partir do Paleozoico.
\end{solution}

\begin{solution}{exo:g7:geological-time:8}
(a) $4540 \div 100 = 45.4$ milhões de anos por metro. (b) $538.8 \div 45.4
\approx 11.9$: cerca de \qty{11.9}{m} antes da linha de chegada. (c) $66 \div
45.4 \approx 1.45$: cerca de \qty{1.45}{m}. (d) $0.3 \div 45.4 \approx
0.0066$~m, ou seja, cerca de \qty{0.66}{cm}, menos de um centímetro: um
corredor mal perceberia.
\end{solution}

\begin{solution}{exo:g7:geological-time:9}
(a) O Siluriano veio primeiro. Sim: a \omterm{def:g3:stones-pebbles-sand:rock}{rocha} mais velha fica embaixo da
mais nova. (b) As camadas cinzentas são depositadas na horizontal; elas
são empurradas até ficarem em pé; elas são desgastadas; a \omterm{def:g3:stones-pebbles-sand:sand}{areia} vermelha
é depositada por cima. (c) É uma superfície de \omterm{def:g5:water-shapes-the-land:erosion}{erosão} entre \omterm{def:g3:stones-pebbles-sand:rock}{rochas} de
idades muito diferentes: o longo tempo durante o qual as camadas
cinzentas foram levantadas e desgastadas não deixou nenhuma camada.
\end{solution}

\begin{solution}{exo:g7:geological-time:10}
(a) Os dois folhelhos com amonites têm a mesma idade, e os dois calcários
com ouriços-do-mar também: mesmos \omterm{def:g7:sediments-and-fossils:fossil}{fósseis}, mesma idade (\omterm{def:g7:geological-time:fossil-succession}{princípio da
sucessão dos fósseis}). (b) A mais velha é o arenito com trilobitas do
penhasco A (ele fica embaixo do folhelho com amonites, que tem a mesma
idade da camada de baixo do penhasco B). A mais nova é a argila com
folhas do penhasco B. (c) De baixo para cima: arenito com trilobitas;
folhelho com amonites; calcário com ouriços-do-mar; argila com folhas.
\end{solution}

\begin{solution}{exo:g7:geological-time:11}
(a) O xisto. Ele fica embaixo do conglomerado (superposição), e os \omterm{def:g3:stones-pebbles-sand:pebble}{seixos}
do conglomerado são pedaços de xisto (inclusão); a superfície de \omterm{def:g5:water-shapes-the-land:erosion}{erosão}
também corta o xisto. (b) Depois: o arenito foi depositado na horizontal,
então ele foi inclinado mais tarde. (c) Antes: a argila cobre a \omterm{def:g7:earthquakes:fault}{falha} sem
ser quebrada. (d) O xisto se forma; ele é desgastado (superfície
ondulada); o conglomerado, o arenito e o calcário são depositados, nessa
ordem; as camadas são inclinadas, e a \omterm{def:g7:earthquakes:fault}{falha} as quebra --- só o corte não
diz qual dessas duas coisas veio primeiro; as camadas inclinadas e a
\omterm{def:g7:earthquakes:fault}{falha} são aplainadas pela \omterm{def:g5:water-shapes-the-land:erosion}{erosão}; a argila é depositada por cima.
(e) Duas: a superfície ondulada entre o xisto e o conglomerado, e a
superfície entre as camadas inclinadas e a argila horizontal.
\end{solution}

\begin{solution}{exo:g7:geological-time:12}
(a) A metade de 365 dias é 182.5 dias. De janeiro a junho são
$31+28+31+30+31+30 = 181$ dias, então a história está pela metade 1.5 dia
depois do começo de julho: em \textbf{2~de~julho, ao meio-dia}. (b) $4540 \div 2 = 2270$ milhões
de anos. (c) $80 \div \num{4540000000} \times \num{31536000} \approx 0.56$:
cerca de \textbf{0.6~segundo}.
\end{solution}

\begin{solution}{pb:g7:geological-time:1}
\textbf{1.} \num{4540} milhões de anos.

\textbf{2.} $4540 \div 365 \approx 12.4$ milhões de anos.

\textbf{3.} $\num{4540000000} \div (365 \times 24) = \num{4540000000}
\div 8760 \approx \num{518000}$ anos (com o valor arredondado de
\num{12.4} milhões, $\num{12400000} \div 24 \approx \num{517000}$).

\textbf{4.} $\num{4540000000} \div \num{525600} \approx 8638$: cerca de
\num{8600} anos.

\textbf{5.} $4540 - 4030 = 510$ milhões de anos depois da formação da
Terra; $510 \div 12.4 \approx 41$ dias. Janeiro tem 31 dias, então 41
dias se passaram no começo de \textbf{11~de~fevereiro}.

\textbf{6.} $4540 - 3500 = 1040$ milhões de anos; $1040 \div 12.4 \approx
83.9$ dias (83.6 com o valor não arredondado $12.44$). Janeiro e fevereiro
somam 59 dias: são 24.6 a 24.9 dias depois do início de março, em \textbf{25~de~março}.

\textbf{7.} $66 \div 12.4 \approx 5.3$ dias antes do fim do ano, ou seja,
$365 - 5.3 = 359.7$ dias depois de 1º~de~janeiro. Os onze primeiros meses
somam 334 dias: 25.7 dias de dezembro, em \textbf{26~de~dezembro}.

\textbf{8.} $538.8 \div 4540 \approx 0.119$: sobram \textbf{11.9~\%} do
ano, ou seja, $0.119 \times 365 \approx 43$ dias. $365 - 43 =
322$ dias depois de 1º~de~janeiro: ele começa em \textbf{novembro} (em
18~de~novembro).

\textbf{9.} $\num{10000} \div 8600 \approx 1.2$ minuto.

\textbf{10.} $5000 \div 8600 \approx 0.58$ minuto, e $0.58 \times 60
\approx 35$: cerca de \textbf{35 segundos} antes da meia-noite.

\textbf{11.} $\num{300000} \div 8600 \approx 34.9$: cerca de \textbf{35
minutos} antes da meia-noite.

\textbf{12.} 35 minutos antes da meia-noite de 31~de~dezembro: a nossa
espécie aparece por volta das \textbf{23:25}, na última hora do ano.
\end{solution}
